\chapter{Measuring What Matters}
\label{ch:measurement}

The last chapter of a book on algorithms should be about how to tell whether one is
faster than another, because that turns out to be harder than either the mathematics or
the implementation.

The chapter has two halves. The first is methodological: which regime you are in
decides what to optimise, and most benchmark tables in this literature do not say which
regime they measured. The second is a worked case study---the long-only portfolio, end
to end---which puts every part of the book together and reports the numbers honestly,
including the ones that make the method look less impressive than a headline figure
would.

\section{Two regimes}

\begin{quote}
Below a few hundred variables, a solve costs what its array operations cost to
\emph{dispatch}, not what its arithmetic costs. Above a few thousand, the reverse.
\end{quote}

This is not a subtlety; it inverts the design objective.

In the dispatch-bound regime, halving the flop count buys nothing and halving the
\emph{iteration} count buys everything. A method that performs more arithmetic in fewer
steps wins. This is why the guessing strategy of Chapter~\ref{ch:guessing} beats the
exact walk of Chapter~\ref{ch:walking} at moderate $n$ by \emph{more} than the gap
between two implementations of that walk in different languages---a striking measurement,
because it says the algorithmic choice dominates the language choice, which is the
opposite of the usual folklore.

In the arithmetic-bound regime, the opposite: iteration counts matter only through their
product with per-step cost, and structure in the operator (Chapter~\ref{ch:operator})
is what pays.

\emph{Knowing which regime you are in is the whole of the decision.} Effort spent on
kernels is wasted where the bottleneck is the interpreter, and effort spent on iteration
counts is wasted where it is not. So the first thing a benchmark should report is not a
speedup but a regime---and the cheapest diagnostic is to plot the local exponent
$d\log t/d\log n$ and look for where the curve leaves its flat overhead floor.

\begin{remark}[The overhead floor is not an asymptotic]
Remark~\ref{rem:exponent} of Chapter~\ref{ch:warmstart} made this point about a specific
measurement and it is worth generalising. A timing curve that is flat at small $n$ is
not exhibiting a better complexity there; it is exhibiting a constant. Fitting a single
power law across both regimes produces an exponent that describes neither. Fit them
separately, and say where you cut.
\end{remark}

\section{Designing a benchmark family that can distinguish}

Two failure modes recur, and both are avoidable.

\paragraph{The family that hides the exponent.} If two candidate complexities coincide
on your test family, no amount of data distinguishes them. On a box-constrained family
the active set grows with $n$, so $O(nk)$ and $O(n^2)$ predict the same curve; holding
$k$ fixed while $n$ grows separates them. Before running a scaling study, write down the
two hypotheses and check that the family can tell them apart.

\paragraph{The family that hides the hazard.} Chapter~\ref{ch:guessing} gave the sharp
case. Over six hundred random instances across four well-posed constraint families, the
rank failure that destroys the finite-termination guarantee \emph{never occurs};
duplicating columns of the constraint matrix raises it to $93\%$. A literature testing
only on the first kind of family would reasonably conclude the hazard is theoretical.

The general point is uncomfortable and worth stating plainly: \emph{test families in a
literature tend to share the structure that makes the literature's methods look good.}
Finding out which structure that is, is part of reading a paper. The synthetic families
common in the QP literature produce small active sets; a long-only portfolio produces a
huge one---on S\&P~500 data roughly $450$ of $495$ constraints active. Those exercise
entirely different code paths, and a method tuned on one may be untested on the other.

\section{What to compare against}

Three levels of comparison give three different numbers, and conflating them is the most
common overstatement in applied optimization papers.

\begin{enumerate}
\item \textbf{Same family, different linear algebra.} The fair algorithmic comparison.
      A dual active-set method on $(J,R)$ factors against a dual active-set method on
      recursive $LDL\T$ updates \cite{arnstrom2022daqp} isolates the linear algebra and
      nothing else.
\item \textbf{Different algorithm, native API.} A Krylov active-set method against an
      interior-point solver \cite{clarabel2024} or an operator-splitting solver
      \cite{osqp}, each called through its own interface. This measures algorithms, and
      it is the number a paper's abstract should carry.
\item \textbf{Different algorithm, through a modelling layer.} The same comparison with
      the competitor wrapped in a general-purpose modelling front end. This measures
      algorithm \emph{plus} problem-construction overhead, and it is an \emph{end-to-end}
      figure, not an algorithmic one.
\end{enumerate}

The gap between levels 2 and 3 is large and systematic. On the portfolio study below,
the algorithmic gain over an interior-point solver is roughly $5.7\times$ without
shrinkage and $22\times$ under heavy shrinkage; the corresponding end-to-end figures
against the same solver behind a modelling layer are $156\times$ and $661\times$. Both
sets of numbers are true. Only the first pair is about the algorithm, and a paper that
reports only the second pair is not lying but is not informing either.

\emph{Report the fair comparison in the abstract and the end-to-end one in the text,
labelled as such.}

\section{Case study: the long-only portfolio, end to end}

Everything in this book meets here.

\paragraph{The problem.} Minimise $w\T\hat\Sig w - \rho\hat\mu\T w$ subject to
$Bw = c$ and $w \ge 0$, with $\hat\Sig$ estimated from a $T \times n$ matrix of
demeaned returns. The data: $494$ S\&P~500 constituents over five years, $T = 1213$,
so $n/T \approx 0.41$.

\paragraph{Step 1: do not form the covariance.} The variance term is
$\norm{Xw}^2/T$, so the gradient is $w \mapsto X\T(Xw)$: two products, $O(nT)$, no
$n \times n$ storage (Chapter~\ref{ch:operator}, Gram backend). The covariance matrix is
not a primitive of the problem; it is a computational detour.

\paragraph{Step 2: eliminate the balance multipliers.} The $p$ multipliers of $Bw = c$
come out analytically through a $p\times p$ Schur complement, leaving $p+1$ SPD solves
(Chapter~\ref{ch:krylov}, Section~3). For the budget constraint alone this is
$w = v_1/(\onevec\T v_1)$ with $v_1 = \hat\Sig^{-1}\onevec$: a \emph{single} solve, and a
formula visible in Merton \cite{merton1972} sixty years ago but not, until recently,
exploited computationally.

\paragraph{Step 3: enforce non-negativity by an active-set loop.} The primal--dual loop
of Chapter~\ref{ch:guessing} wraps the solver unchanged, with the finite-termination
guarantee of Theorem~\ref{thm:termination} and the inexactness transfer of
Lemma~\ref{lem:inexact}.

\paragraph{Step 4: condition it.} Ledoit--Wolf shrinkage at the intensity practitioners
already use, or RMT cleaning (Chapter~\ref{ch:conditioning}).

\paragraph{Step 5: for a frontier, warm start.} Chapter~\ref{ch:warmstart}.

\subsection{What is measured}

\begin{center}
\footnotesize
\begin{tabular}{@{}lll@{}}
\toprule
Quantity & Without shrinkage & Heavy shrinkage ($\alpha = 0.5$) \\
\midrule
$\kappa(\hat\Sig)$                    & $\approx 109{,}000$ & $\approx 147$ \\
CG iterations                         & $684$   & $82$ \\
Outer active-set steps                & $18$    & $6$ \\
Speedup vs.\ interior point (native)  & $5.7\times$ & $22\times$ \\
Speedup vs.\ interior point (modelling layer) & $156\times$ & $661\times$ \\
\bottomrule
\end{tabular}
\end{center}

Four readings, in decreasing order of how flattering they are.

\emph{The conditioning effect is real and compound.} Shrinkage cuts CG iterations by a
factor of eight, as Proposition~\ref{prop:cg} predicts from the $\kappa$ collapse. It
also cuts the \emph{outer} count from $18$ to $6$, which is not a $\kappa$ effect at
all: a lower condition number moves the unconstrained portfolio toward equal weights, so
fewer assets receive negative weights in the first primal step. Two independent
mechanisms, one parameter.

\emph{Scaling favours the method.} At $n = 3000$ on synthetic factor-model data with $T$
fixed, the matrix-free method takes about $0.081$ seconds and its runtime grows
substantially \emph{sub}quadratically, while interior-point methods scale as $O(n^3)$
per iteration. The gap widens with $n$.

\emph{A well-engineered general solver is still slower, and this is the informative
comparison.} An operator-splitting solver called through its \emph{direct} API---no
modelling overhead---remains $6$ to $10\times$ slower. That is the honest statement of
what exploiting structure buys: not three orders of magnitude, but roughly one, and
reliably.

\emph{The three-order-of-magnitude figures are mostly about something else.} A $50$-point
frontier for $n = 500$ completes about $1{,}455\times$ faster than a naive loop over a
modelling layer. That number is real and it is what a user experiences, but most of it
is problem-construction overhead paid $50$ times, not algorithmic superiority. Say so.

\subsection{The statistical objective, checked separately}

A faster answer to the wrong problem is worthless, so the conditioning choice has to be
validated on its own terms. A rolling-window out-of-sample study on the same universe
finds realised out-of-sample variance minimised at $\alpha \approx 0.3$---inside the
heavy-shrinkage range that gives the conditioning benefit---with turnover falling
monotonically in $\alpha$. And RMT-cleaned estimates carry no measurable statistical
penalty relative to standard shrinkage on these data.

So the statistical and numerical objectives are reconciled at the intensities
practitioners already use, and the apparent tension noted in
Chapter~\ref{ch:conditioning} is an artefact of the Frobenius surrogate.

\emph{Notice what this required: a separate experiment, with a separate metric, that
could have come out the other way.} A computational paper that only measures time has
not established that its method should be used.

\section{Reporting}

A checklist for a benchmark section, derived from the above.

\begin{enumerate}
\item State the regime. Report where the overhead floor ends.
\item Report the \emph{fair} comparison---same family, or native API---as the headline,
      and label end-to-end figures as end-to-end.
\item Report iteration counts and wall time separately, and never mix warm and cold
      iteration counts in one column (Chapter~\ref{ch:warmstart}).
\item Say what the test family's active-set size is, since methods differ enormously
      between small and large active sets.
\item Design the family so that competing complexity hypotheses predict different
      curves, and say which hypotheses you were testing.
\item Verify correctness independently of timing: every returned point should pass the
      certificate of Chapter~\ref{ch:certificate} against the original data, and agree
      with a reference solver.
\item Where a method wins by an approximation---a cleaned estimator, a regularised
      operator---validate that approximation on its own metric.
\item Report what did not work, and on what family the hazard you did not hit would have
      fired.
\end{enumerate}

\section{Closing}

The book began with the claim that the analysis of the convex quadratic program is easy
and the combinatorics hard, and that every method is a strategy for searching among
active sets. Four strategies later, the strategies have turned out to share more than
they differ: one certificate makes all of them safe; one operator interface serves all
of them; one linear-algebra kernel---a free block, bordered by active normals, updated
rather than rebuilt---sits inside all of them; and one parameter, allowed to vary, turns
the whole subject into a single curve that four literatures traced independently.

If there is one thing to carry away, it is the certificate principle of
Chapter~\ref{ch:certificate}, because it generalises far past quadratic programming: a
method that may fail, but that can recognise its own success, needs a certificate and
not a convergence theory. Ask of any algorithm you meet whether it can check its own
answer. If it can, you may be as inventive as you like about how it finds one.

\begin{notes}
The regime observation and the measurement that the fast path beats the exact walk by
more than the gap between two language implementations are from the Goldfarb--Idnani note \cite{schmelzer2026quadprog}. The
portfolio measurements are from \cite{schmelzer2026matrixfree}; the frontier-sweep and
RMT figures from \cite{schmelzer2026rmt}; the out-of-sample study from
\cite{schmelzer2026matrixfree}, with \cite{demiguel2009} for the wider evidence on
shrinkage and out-of-sample performance. The solvers benchmarked against are Clarabel
\cite{clarabel2024} and OSQP \cite{osqp}, with CVXPY \cite{diamond2016} the modelling
layer.
\end{notes}

\begin{exercises}

\begin{exercise}
Given timings $t(n)$ at $n = 50, 100, 200, 400, 800, 1600$, describe how to estimate the
local exponent $d\log t/d\log n$ and how to locate the end of the overhead floor.
\end{exercise}

\begin{exercise}
Explain why a single power-law fit across both regimes yields an exponent that describes
neither, and construct synthetic timings for which the fitted exponent is $1.0$ although
neither regime has that exponent.
\end{exercise}

\begin{exercise}
For each of the three comparison levels of Section~3, name a question it answers and a
question it cannot.
\end{exercise}

\begin{exercise}
The outer active-set count fell from $18$ to $6$ under shrinkage. Give the mechanism, and
explain why it is \emph{not} a consequence of Proposition~\ref{prop:cg}.
\end{exercise}

\begin{exercise}
Suppose a method is $1{,}455\times$ faster than a baseline end to end and $22\times$
faster algorithmically. What is the remaining factor measuring, and how would you
measure it directly?
\end{exercise}

\begin{exercise}
Design a test family on which the block-pivoting rank hazard of
Chapter~\ref{ch:guessing} fires, and one on which it never does. What does the existence
of both say about published benchmark suites?
\end{exercise}

\begin{exercise}[Implementation]
Assemble the full pipeline of Section~4 on real or simulated returns: Gram operator,
Schur elimination, active-set loop, shrinkage, warm-started frontier sweep. Verify every
returned portfolio against the certificate of Chapter~\ref{ch:certificate} and against a
general-purpose QP solver. Then reproduce the table of Section~4.1 on your own data and
comment on every figure that differs.
\end{exercise}

\begin{exercise}[Implementation]
Run a rolling-window out-of-sample study over shrinkage intensities
$\alpha \in \{0, 0.017, 0.1, 0.3, 0.5, 0.7\}$, recording realised out-of-sample variance
and turnover at each. Does your data reproduce the reconciliation of Section~4.2? If not,
what differs?
\end{exercise}

\end{exercises}
